\documentclass[11pt]{article}
\usepackage{mathblatt}
\makeatletter
\setlength{\mbpfbe}{0mm}% keine Punkte (Bauregel 6.4)
\makeatother
% Rechenraum als Karo (Bauregel 6.9): n Rechenschritte = n mal 10 mm
\newcommand{\karo}[1]{\par\smallskip\noindent\begin{tikzpicture}
  \clip (0,0) rectangle ({\linewidth-1mm},{#1*10mm});
  \draw[black!16,line width=0.3pt,step=5mm] (0,0) grid ({\linewidth},{#1*10mm});
  \end{tikzpicture}\par}
\newcommand{\linie}[1]{\,{\color{black!45}\rule[-1pt]{#1}{0.4pt}}\,}
\newcommand{\pkt}{\textperiodcentered{}}
\begin{document}

\pfheftstil
\pflaufkopf[P10 \pkt{} mehr]{Quadratische Gleichung lösen}{PA40-3}\pfseitevon
\pfheftkopf{Quadratische Gleichung lösen}{mehr}
\pfstufenzeile{Nullstellen berechnen}
% 2020-OS-K3e-m1 (Schwere II, rang 7)
\begin{pfaufg}{nach~P10\\’20~\pkt{}~3e}{1.}{}
Berechne die Nullstellen der Funktion mit $y = 2x^2 + 6x - 20$.
\karo{3}
\fusshilfe{\mbox{1:~2; $-$5}}
\end{pfaufg}
% 2017-OS-K5d-m1 (Schwere III, rang 8)
\begin{pfaufg}{nach~P10\\’17~\pkt{}~5d}{2.}{}
Die Parabel $p$ hat die Gleichung $y = (x+4)^2 - 3$ und ist im Koordinatensystem gezeichnet. Zeige, dass man die Gleichung von $p$ auch als $y = x^2 + 8x + 13$ schreiben kann. Berechne die Nullstellen von $p$.
\par\smallskip\noindent{\small\itshape Skizze: Graph: p: (x+4)$^{2}$-3; x -8..0; y -4..4; Gitter ja}\par
\karo{4}
\fusshilfe{\mbox{2:~x² + 8x + 13; $-$2,27; $-$5,73}}
\end{pfaufg}
\par\vfill\noindent{\scriptsize\color{mbgrau}Originale: P10 2020 \pkt{} 3e, P10 2017 \pkt{} 5d\par}
\end{document}